\documentclass[a4paper]{article}

\usepackage{cv}
\hypersetup{pdftitle={林彦 - 求职信}}

\begin{document}
\cvname{林彦}
\cvtagline{助理教授 / 计算机科学系 / 奥尔堡大学}

电子邮箱：\href{mailto:lyan@cs.aau.dk}{lyan@cs.aau.dk} \,/\,
个人主页：\href{https://www.yanlincs.com}{yanlincs.com} \,/\,
电话：+45 9186-1833

{\color{rule}\rule{\linewidth}{0.4pt}}

\vspace{\entryskip}
\today

\vspace{\entryskip}
尊敬的评审委员会成员：

\vspace{\entryskip}
您好！
我谨申请[单位][院系]的[职位]。
我叫林彦，现任奥尔堡大学计算机科学系助理教授，2024 年 3 月于北京交通大学获得计算机科学博士学位。
2022 年 9 月，我以访问博士生身份首次加入奥尔堡大学，2024 年 6 月回到奥尔堡大学担任博士后研究员，2026 年 6 月起担任现职。
我期待有机会在这一职位上继续从事学术工作。

\vspace{\entryskip}
我的博士论文《时空轨迹数据的自监督学习》提出面向时空轨迹数据的自监督表示学习方法。
自 2019 年以来，我在 IEEE TKDE、SIGMOD、NeurIPS、KDD、WWW、AAAI 和 IJCAI 等顶级期刊和会议发表同行评议论文 33 篇，其中 17 篇以第一作者、共同第一作者或主要学生作者身份完成，一项工作获选 NeurIPS 2025 口头报告。
论文在谷歌学术上的总被引次数超过 1,000，h 指数为 15。
我的研究聚焦时空数据挖掘，主要开发面向轨迹表示、轨迹恢复和出行时间估计等任务的人工智能模型。

\vspace{\entryskip}
我也参与跨学科的人工智能合作研究。
我参与 Villum Synergy Programme 资助的“基于扩散概率模型的材料逆向设计”项目，与奥尔堡大学化学与生物科学系合作，将人工智能应用于无序材料的发现和设计。
我曾主持科研项目，并参与 Villum Foundation 和中国国家自然科学基金资助项目的申请。
入职后，我计划申请 Villum Synergy、Villum Experiment 和 DFF Research Project 等资助，以支持研究并为学生提供科研训练机会。

\vspace{\entryskip}
我将科研融入教学。
我讲授的“AI Systems \& Infrastructure”课程以我的研究和项目实践为基础，涵盖人工智能模型部署、可扩展基础设施和适用于生产环境的人工智能系统。
我遵循奥尔堡大学的基于问题的学习模式，从零编写通俗易读的博客式课程材料，并指导应用人工智能方向的本科生和硕士生学期项目。
学生对课程结构、内容和教学质量给予积极评价，认为我的指导很有帮助，并认为我平易近人、易于沟通。
入职后，我将开设新课程，并指导连接科研与实际应用的学位论文项目。

\vspace{\entryskip}
我积极服务学术共同体。
我现任 IEEE 计算机学会丹麦分会秘书，担任 IEEE TKDE、ACM TKDD、IEEE TNNLS、ACM TIST、IEEE TII 和 IEEE TVT 等期刊的审稿人，以及 ICLR、NeurIPS、KDD、CVPR、ICCV、AAAI、IJCAI、WWW 和 WACV 等重要会议的程序委员会委员。
入职后，我也愿意承担课程协调和院系委员会工作等行政职责。

\vspace{\entryskip}
随信附上个人简历、完整论文列表和教学档案。
如有需要，以下人士可为我提供推荐信：曾与我开展科研合作的 Christian S. Jensen 教授和 Bin Yang 教授，与我合作开展 Villum Synergy Programme 项目的奥尔堡大学化学与生物科学系 Morten M. Smedskjaer 教授，以及与我密切合作开展课程教学和学生项目指导的 Andreas Aakerberg 博士。
如需其他材料，我很乐意进一步提供。
感谢您的审阅与考虑。

\vspace{\entryskip}
此致\\
敬礼！

\vspace{\entryskip}
林彦

\end{document}
